\section{Rank approximation and projections in the tracial ultrapower}
\label{u:sec:ranks}

Let $A$ be simple, unital and infinite dimensional, with
$T=T(A)\ne\varnothing$.  We first approximate the trace function of a
positive element by the rank of a finite-matrix positive element.
The construction transfers the trace function exactly to a hereditary
subalgebra of small support rank, then takes integer spectral cuts.
The numerical argument requires neither separability, nuclearity nor
comparison.  We will then add strict comparison to place these rank
approximants in prescribed hereditary subalgebras and obtain projections
in the uniform tracial ultrapower.  All matrix traces are unnormalized.

The spectral discretization of traces into dimension functions is related
to the rank approximation argument of D\u ad\u arlat and Toms
\cite[proof of Theorem~5.1]{DadarlatToms}. The trace-preserving compression
below concentrates the rounding error on one small hereditary support;
we prove every step under the displayed elementary hypotheses.

The small-support input is Lemma~\ref{u:lem:thin-full}, which was
proved in Section~\ref{r:sec:projections} under these same elementary
hypotheses.  In the notation of Section~\ref{u:sec:preliminaries}, it
provides, for every $\delta>0$, a norm-one full positive element $e$
with $d_\tau(e)<\delta$
for every $\tau\in T$.

\begin{proposition}[Uniform approximation of trace values by ranks]
\label{u:prop:rank-approximation}
Let $A$ be a simple, unital, infinite-dimensional $C^*$-algebra with
$T=T(A)\ne\varnothing$.  For every $b\in A_+$ and $\delta>0$ there
are an integer $L\ge1$ and $v\in M_L(A)_+$ such that
\begin{equation}
 \tau(b)-\delta<d_\tau(v)\le\tau(b)
 \qquad(\tau\in T).
 \label{u:eq:rank-approximation}
\end{equation}
This statement requires neither separability, nuclearity nor strict
comparison.  There is no asserted bound on $L$ independent of $b$ and
$\delta$.
\end{proposition}

\begin{proof}
Take $e$ from Lemma~\ref{u:lem:thin-full}.  We first extract an exact
finite identity from its fullness.  Since the linear span of $AeA$ is
dense in $A$, choose $u_i,v_i\in A$, $1\le i\le n$, so that
\[
 s=\sum_{i=1}^n u_i e v_i,\qquad \norm{1-s}<1.
\]
Thus $s$ is invertible.  Form the row and column
\[
 U=(u_1e^{1/2},\ldots,u_ne^{1/2}),\qquad
 V=(e^{1/2}v_1,\ldots,e^{1/2}v_n)^{\mathsf T}.
\]
Since $s=UV$, the element $y=V^*V=\sum_i v_i^*ev_i$ satisfies
\[
 \norm{s^{-1}}^{-2}1\le s^*s
       =V^*U^*UV\le\norm{U}^2 y.
\]
In particular $y$ is invertible.  Setting $x_i=v_i y^{-1/2}$ gives
the exact identity
\begin{equation}
 \sum_{i=1}^n x_i^* e x_i=1.
 \label{u:eq:fullness-frame}
\end{equation}

Define
\[
 c=\sum_{i=1}^n e^{1/2}x_i b x_i^*e^{1/2}\in\Her(e)_+.
\]
Additivity and traciality, followed by
Equation~\eqref{u:eq:fullness-frame}, give, for every $\tau\in T$,
\begin{equation}
 \tau(c)=\sum_i\tau(bx_i^*ex_i)=\tau(b),
 \qquad d_\tau(c)\le d_\tau(e)<\delta.
 \label{u:eq:trace-compression}
\end{equation}
The norm of $c$ is finite but may be large.  Choose an integer
$L\ge\max\{1,\norm{c}\}$ and put
\[
 v=\bigoplus_{j=1}^{L}(c-j)_+\in M_L(A)_+.
\]
For $0\le s\le\norm{c}$, let
$N_L(s)=\sum_{j=1}^L\mathbf 1_{(j,\infty)}(s)$.  The elementary
counting identity gives
\begin{equation}
 0\le s-N_L(s)\le\mathbf 1_{(0,\infty)}(s).
 \label{u:eq:integer-cuts}
\end{equation}
For nonintegral positive $s$ the count is $\lfloor s\rfloor$; at a
positive integer $s$ it is $s-1$, so the difference is exactly one.
At $s=0$ the difference is zero.  This also checks the upper endpoint
when $\norm{c}=L$.

Apply bounded Borel functional calculus and the normal extension of
$\tau$ to Equation~\eqref{u:eq:integer-cuts}.  The unnormalized matrix
trace adds the support ranks of the direct summands, and hence
\[
 d_\tau(v)=\sum_{j=1}^L\tau\bigl(\mathbf 1_{(j,\infty)}(c)\bigr),
 \qquad
 0\le\tau(c)-d_\tau(v)\le d_\tau(c)<\delta.
\]
Together with Equation~\eqref{u:eq:trace-compression}, this is
Equation~\eqref{u:eq:rank-approximation}.
\end{proof}


We now assume in addition that $A$ has strict comparison as in
Definition~\ref{u:def:strict-comparison}, and retain its uniform tracial
ultrapower $D$ and limit traces $\Lambda$ from
Section~\ref{u:sec:preliminaries}.  Comparison will first place the rank
approximants inside a prescribed hereditary subalgebra of $A$.  Passing
to $D$ will turn them into projections with exact spectral support.

We use the support and normal-trace conventions of
Section~\ref{u:sec:preliminaries}.  In particular, for $a\in M_n(A)_+$,
\[
 d_\tau(a)=(\tau\otimes\Tr)(\supp(a)),
 \qquad \supp(a)=\mathbf 1_{(0,\infty)}(a)\in M_n(A)^{**}.
\]
The right side uses the normal extension of the matrix trace.  Every
element of $\Her(a)$ is supported on both sides by $\supp(a)$.

\begin{lemma}[Almost projections in a prescribed hereditary subalgebra]
\label{u:lem:almost-projection}
For $b\in A_+$ with $\norm{b}<1$ and $\delta>0$, there is a positive
contraction $p_0\in\Her(b)$ such that, for every $\tau\in T$,
\begin{equation}
 \tau(b)-2\delta<\tau(p_0)\le\tau(b),
 \qquad 0\le\tau(p_0-p_0^2)<2\delta.
 \label{u:eq:almost-projection}
\end{equation}
\end{lemma}

\begin{proof}
If $b=0$, take $p_0=0$.  Otherwise choose $v\in M_L(A)_+$ as in
Proposition~\ref{u:prop:rank-approximation}.  Faithfulness of each trace
and the support formula give
\[
 0<d_\tau(b),\qquad
 d_\tau(v)\le\tau(b)\le\norm{b}\,d_\tau(b)<d_\tau(b).
\]
Strict comparison, applied in $A\otimes\K$, gives $v\precsim b$.

There is a single $t>0$ such that
\begin{equation}
 d_\tau((v-3t)_+)>\tau(b)-2\delta
 \qquad(\tau\in T).
 \label{u:eq:uniform-rank-cut}
\end{equation}
Here is the uniformity argument.  For fixed $t>0$, the rank on the
left is a lower semicontinuous function of $\tau$: it is the supremum
over $m\in\N$ of the continuous functions
\[
 \tau\longmapsto(\tau\otimes\Tr)
       \bigl(\min\{1,m(v-3t)_+\}\bigr).
\]
Moreover, these cut ranks increase to $d_\tau(v)$ as $t$ decreases to
zero.  Thus the open sets
\[
 U_t=\{\tau\in T:
       d_\tau((v-3t)_+)>\tau(b)-2\delta\}
\]
increase as $t$ decreases and cover $T$, by
Equation~\eqref{u:eq:rank-approximation}.  Compactness of $T$ supplies
a finite subcover.  The smallest of its cut parameters gives
Equation~\eqref{u:eq:uniform-rank-cut}.

We next obtain a comparison witness of the required matrix shape.
Regard $b$ as supported in the first matrix corner of $A\otimes\K$,
and $v$ as supported in its first $L$ corners.  Let $P_1$ and $P_L$
denote the corresponding corner projections.  From $v\precsim b$,
choose $X\in A\otimes\K$ with $\norm{X^*bX-v}<t$.  The row
$x=P_1XP_L\in M_{1,L}(A)$ satisfies
\[
 w=x^*bx=P_LX^*bXP_L,\qquad \norm{w-v}<t.
\]
Set $\alpha=\norm{w-v}$ and $\gamma=t-\alpha>0$.  In
$M_L(A)^{**}$ put
$E=\mathbf 1_{(3t,\infty)}(v)$ and
$F=\mathbf 1_{(t,\infty)}(w)$.  The order inequalities
$(w-2t)_+\ge w-2t1$ and $w\ge v-\alpha1$ give
\[
 E(w-2t)_+E
 \ge E(v-(2t+\alpha)1)E\ge\gamma E.
\]
Similarly,
\[
 FvF\ge F(w-\alpha1)F\ge\gamma F.
\]
For a positive element $a$ in a von Neumann algebra and a projection
$Q$ with $QaQ\ge\gamma Q$, $\gamma>0$, the polar decomposition of
$a^{1/2}Q$ has initial projection $Q$ and final projection at most
$\supp(a)$.  Indeed its absolute square is $QaQ$, which is invertible
in the $Q$-corner, while its range is contained in the support of $a$.
Applying this observation to the preceding two inequalities and then
applying each normal trace gives
\begin{equation}
 d_\tau((v-3t)_+)\le d_\tau((w-2t)_+),
 \qquad d_\tau((w-t)_+)\le d_\tau(v).
 \label{u:eq:perturbed-ranks}
\end{equation}
The positive constant $\gamma$ is independent of $\tau$.

Set $a=b^{1/2}x\in M_{1,L}(A)$ and
$w'=aa^*=b^{1/2}xx^*b^{1/2}\in\Her(b)_+$.  In the bidual of the
linking matrix algebra, the polar decomposition $a=u|a|$ identifies
the positive spectral projections of $a^*a=w$ and $aa^*=w'$:
\[
 \mathbf 1_{(s,\infty)}(w')
    =u\mathbf 1_{(s,\infty)}(w)u^*
 \qquad(s>0).
\]
Consequently $w$ and $w'$ have the same cut ranks, with the
unnormalized matrix trace on the $w$ corner and $\tau$ on the $w'$
corner.

Choose continuous $g:[0,\infty)\to[0,1]$ which is zero on $[0,t]$
and one on $[2t,\infty)$, and put $p_0=g(w')$.  This is a positive
contraction in $\Her(b)$ since $g(0)=0$.  Scalar functional calculus
gives
\[
 \mathbf 1_{(2t,\infty)}\le g\le\mathbf 1_{(t,\infty)},
 \qquad
 0\le g-g^2\le
       \mathbf 1_{(t,\infty)}-\mathbf 1_{(2t,\infty)}.
\]
Taking traces and using the equality of cut ranks yields
\[
 \begin{split}
 d_\tau((w-2t)_+)&\le\tau(p_0)\le d_\tau((w-t)_+),\\
 \tau(p_0-p_0^2)&\le
       d_\tau((w-t)_+)-d_\tau((w-2t)_+).
 \end{split}
\]
Equations~\eqref{u:eq:rank-approximation},
\eqref{u:eq:uniform-rank-cut} and \eqref{u:eq:perturbed-ranks} now imply
Equation~\eqref{u:eq:almost-projection}.  All ranks here are finite,
so the displayed differences are well defined.  The argument also
allows $v=0$; its rank approximation then gives $\tau(b)<\delta$
at every trace, and $x=0$ is a permissible witness.
\end{proof}

\begin{proposition}[Projection supply with exact spectral support]
\label{u:prop:projection-supply}
For every positive contraction $h\in D$ and $0<\eta<1$, there is a
projection $p\in D$ such that
\begin{equation}
 p\le\mathbf 1_{[\eta,1]}(h)\quad\text{in }D^{**},
 \qquad
 \tau(p)=\tau((h-\eta)_+)
       \quad(\tau\in\Lambda).
 \label{u:eq:projection-supply}
\end{equation}
In addition, for $k_\eta(s)=\min\{s/\eta,1\}$ on $[0,1]$, one has
\begin{equation}
 k_\eta(h)p=p=pk_\eta(h)\quad\text{in }D.
 \label{u:eq:projection-cutoff}
\end{equation}
\end{proposition}

\begin{proof}
Choose positive contractive lifts $(h_\ell)_{\ell\in\N}$ of $h$.
For $b_\ell=(h_\ell-\eta)_+$ we have
$\norm{b_\ell}\le1-\eta<1$.  Apply
Lemma~\ref{u:lem:almost-projection} with $\delta=1/\ell$ to obtain
positive contractions $p_\ell\in\Her(b_\ell)$ satisfying
\[
 \sup_{\tau\in T}|\tau(p_\ell)-\tau(b_\ell)|\le2/\ell,
 \qquad
 \sup_{\tau\in T}\tau(p_\ell-p_\ell^2)\le2/\ell.
\]
Let $p$ be the class of $(p_\ell)$ in $D$.  Since $p_\ell$ is a
positive contraction,
\[
 0\le(p_\ell-p_\ell^2)^2\le p_\ell-p_\ell^2.
\]
Its projection defect therefore has uniform tracial $2$-norm tending
to zero, so $p=p^*=p^2$ in $D$.  For an arbitrary limit trace
$\tau=\lim_\omega\tau_\ell$ on represented sequences, the first
estimate gives
\[
 \tau(p)=\lim_{\ell\to\omega}\tau_\ell(b_\ell)
        =\tau((h-\eta)_+).
\]

For the support assertion, functional calculus in $A^{**}$ gives
\[
 \supp(b_\ell)=\mathbf 1_{(\eta,1]}(h_\ell),
 \qquad k_\eta(h_\ell)\supp(b_\ell)=\supp(b_\ell).
\]
Every element of $\Her(b_\ell)$ is supported on both sides by this
projection.  Hence $k_\eta(h_\ell)p_\ell=p_\ell$ exactly at every
coordinate.  Passing to the quotient and taking adjoints proves
Equation~\eqref{u:eq:projection-cutoff}.

Finally represent $D^{**}$ faithfully on a Hilbert space.  The identity
$(1-k_\eta(h))p=0$ says that the range of $p$ lies in the kernel of
$1-k_\eta(h)$.  Its kernel projection is
\[
 \mathbf 1_{\{1\}}(k_\eta(h))
       =\mathbf 1_{[\eta,1]}(h).
\]
Thus $p$ has the asserted order bound in $D^{**}$.  This argument uses
an exact identity in $D$ and does not require any normal trace on
$D^{**}$ to be faithful.
\end{proof}

\begin{corollary}[Trace-zero symmetries in every projection corner]
\label{u:cor:corner-symmetry}
For every projection $q\in D$ there is $r=r^*\in qDq$ such that
\[
 r^2=q,\qquad \tau(r)=0\quad(\tau\in\Lambda).
\]
\end{corollary}

\begin{proof}
Apply Proposition~\ref{u:prop:projection-supply} to $h=q$ and
$\eta=1/2$.  It gives a projection $p\le q$ with
$\tau(p)=\tau(q)/2$ for all $\tau\in\Lambda$.  Then $r=2p-q$
belongs to $qDq$, is selfadjoint, and satisfies
$r^2=4p-4p+q=q$ and $\tau(r)=0$.  This includes $q=0$.
\end{proof}
